% TOP_PROBLEM: {"id": 5200017, "title": "Open Problems on Billiards and Geometric Optics", "category_id": 11, "category": {"id": 11, "name": "dynamical_systems", "display_name": "Dynamical Systems", "description": "Problems about long-term behavior of deterministic systems, Hamiltonian dynamics, and periodic orbits.", "slug": "dynamical-systems", "order_index": 11, "created_at": "2026-07-31T15:26:25.671Z"}, "status": "open"}
% FIRST_POSTED: null
% ENTRY_KIND: statement_only
% Imported from: batch03.tex; UnsolvedMath ID 5200017
\documentclass[11pt]{article}
\usepackage[margin=1in]{geometry}
\usepackage{amsmath,amssymb,mathtools}
\usepackage{fontspec,unicode-math}
\setmainfont{Latin Modern Roman}
\setsansfont{Latin Modern Sans}
\setmathfont{Latin Modern Math}
\usepackage{xurl,hyperref}
\hypersetup{colorlinks=true,urlcolor=blue,linkcolor=blue}
\newcommand{\sref}[2]{\hyperlink{source:#1:#2}{[S#2]}}
\newcommand{\cataloguescope}[1]{\paragraph{Recorded mathematical question.}#1}
\begin{document}
% UnsolvedMath ID 5200017; source catalogue code AMR-051-0017
\setcounter{section}{5200016}
\section{Open Problems on Billiards and Geometric Optics}\label{5200017}
{\small\sffamily UnsolvedMath ID 5200017 · Dynamical Systems}
\subsection{Definitions and mathematical statement}

\paragraph{Shared notation.}$\mathbb N=\{1,2,\ldots\}$, and $\mathbb Z,\mathbb Q,\mathbb R,\mathbb C$ denote the integers, rationals, reals and complex numbers. Any other conventions are specified locally.


An oval is a smooth closed strictly convex plane curve. Its outer billiard reflects an exterior point in the midpoint of a tangent segment, with a consistent tangent branch. A periodic outer orbit determines a circumscribed polygon of extremal area. The area spectrum is the set of these polygonal areas. The question is a spectral research programme and does not specify an operator in advance.
\paragraph{Statement recorded in the supplied source.}
For an oval $\gamma$, the area spectrum of its outer billiard is the set of areas of the circumscribed polygons formed by periodic outer-billiard trajectories. Is this area spectrum related to the spectrum of a differential operator? \sref{5200017}{2}
\subsection{Short English statement}
Find whether the areas of periodic outer-billiard polygons can be related to the spectrum of a differential operator.
\subsection{Sources}
\begin{description}
\item[\hypertarget{source:5200017:1}{S1}] ULAM.AI and the UnsolvedMath Contributors, \emph{UnsolvedMath: A Curated Collection of Open Mathematics Problems}, record 5200017 (AMR-051-0017). CC BY 4.0. \url{https://huggingface.co/datasets/ulamai/UnsolvedMath}.
\item[\hypertarget{source:5200017:2}{S2}] M. Bialy, C. Fierobe, A. Glutsyuk, M. Levi, A. Plakhov and S. Tabachnikov, \emph{Open problems on billiards and geometric optics} (2021), arXiv:2110.10750v2, Serge Tabachnikov, Problem 5, pp. 12--13. \url{https://arxiv.org/pdf/2110.10750v2}.
\end{description}
\label{end:5200017}
\end{document}
